\documentclass[12pt,a4paper]{article}
\usepackage{ugmath}
\usepackage{float}
\usepackage{placeins}
\usepackage{array}
\usepackage[labelfont=bf,labelsep=space]{caption}
\newcommand{\paren}[1]{\left( #1 \right)}
\newcommand{\alumno}{Ricardo León Martínez}
\newcommand{\materia}{Algebra Lineal I}
\newcommand{\profesor}{Claudia Reynoso Alcántara}
\newcommand{\tarea}{Tarea 6}
\newcommand{\fecha}{9/3/2026}

\begin{document}

    \begin{center}
        {\large\textbf{UNIVERSIDAD DE GUANAJUATO}}\\[0.3cm]
        {\normalsize\textbf{DIVISIÓN DE CIENCIAS NATURALES Y EXACTAS}}\\
        {\normalsize\textbf{CAMPUS GUANAJUATO}}\\[1cm]

        {\Large\textbf{\tarea\ (\materia)}}\\[1cm]
    \end{center}

    \noindent
    \textbf{Nombre:} \alumno 
    \hfill 
    \textbf{Fecha:} \fecha 
    \hfill 
    \textbf{Calificación:} \rule{3cm}{0.4pt}

    \vspace{0.3cm}

    \begin{mdframed}[style=mdbluebox,frametitle={Ejercicio 1}]
        Sea $V$ un espacio vectorial sobre un campo $F$ y sean
        $W_{1},\dots,W_{k}$ subespacios vectoriales de $V$.
        Demuestra que
        \begin{equation*}
            W_{1}+\cdots+W_{k}:=\{w_{1}+\cdots+w_{k}:w_{i}\in W_{i},\forall i\in\{1,\dots,n\},
        \end{equation*}
        es un subespacio vectorial de $V$.
    \end{mdframed}

    \begin{mdframed}[style=mdbluebox,frametitle={EJercicio 2}]
        Falso o verdadero (justifica tus respuestas: si es verdadero
        demuestralo y si es falso da un contraejemplo):
        \begin{enumerate}
            \item[(a)] El vacío es un subconjunto linealmente independiente.
            \item[(b)] Sea $V$ un espacio vectorial sobre $F$, si $S\subset V$
            es linealmente dependiente entonces cualquier subconjunto propio
            de $S$ es linealmente independiente.
            \item[(c)] Sea $V$ un espacio vectorial sobre $F$, si $S\subset V$
            es linealmente independiente entonces cualquier subconjunto
            de $S$ es linealmente independiente.
            \item[(d)] Si un conjunto de 4 vectores en $\mathbb{R}^{4}$ es
            linealmente independiente, entonces necesariamente genera a $\mathbb{R}^{4}$
        \end{enumerate}
    \end{mdframed}

    \begin{mdframed}[style=mdbluebox,frametitle={Ejercicio 3}]
        Demuestra que en el espacio vectorial $V=\{f:\mathbb{R}\to\mathbb{R}:f\text{ es función}\}$
        sobre $\mathbb{R}$, existe un conjunto infinito que es linealmente independiente.
    \end{mdframed}

    \begin{mdframed}[style=mdbluebox,frametitle={Ejercicio 4}]
        Encuentra un conjunto $\{v_{1},v_{2},v_{3}\}$ de vectores linealmente dependientes en
        $\mathbb{R}^{3}$ tales que dos a dos son linealmente independientes.
    \end{mdframed}

    \begin{mdframed}[style=mdbluebox,frametitle={Ejercicio 5}]
        Sean $a,b,c\in\mathbb{R}$ distintos entre sí. Demuestra que
        $\{(1,1,1),(a,b,c),(a^{2},b^{2},c^{2})\}$ es linealmente independiente
        en $\mathbb{R}^{3}$.
    \end{mdframed}
\end{document}